% ===================================================================
%  GAMBAR No. 13 — Hotel. --- Restoran. --- Warung Kopi.
%
%  Traduction, faite en 2026, en javanais d'aujourd'hui, sur l'ido de
%  texto/io. Regles de la colonne : voir 10-gambar-01.tex. Une seule
%  numerotation.
%
%  LE MENU (9) EST LE SEUL ENDROIT DU LIVRET OU LA COLONNE DOIT
%  NOMMER DES PLATS QU'AUCUN JAVANAIS N'A MANGES. Les tripes a la
%  mode de Caen, les crevettes au beurre, le gigot de mouton, les
%  epinards a la creme, le Saint-Emilion : ce sont des noms de
%  choses, pas de mots, et la colonne les garde tels quels ou les
%  decrit. « babat cara Caen » nomme les tripes par leur mot javanais
%  et laisse la ville en francais, parce que c'est la recette qui est
%  nommee, non l'organe. Le beurre est « mentega », qui est
%  portugais ; le fromage « keju », portugais aussi. Les deux sont
%  arrives par les memes navires que meja et jendela, et personne ne
%  les sent etrangers depuis quatre cents ans.
%
%  ET C'EST LE TABLEAU OU LE MOT « warung » NE PEUT PAS SERVIR, ET
%  C'EST DOMMAGE. Le cafe (73, 74) est un etablissement a billard,
%  echecs, dames, dominos et trictrac ; le warung javanais est une
%  echoppe de rue a un ou deux bancs. « warung kopi » est ecrit dans
%  le titre parce que l'ido dit « kafeerio » et que le javanais n'a
%  rien d'autre — mais le lieu grave est un cafe europeen, et la
%  colonne ne le transpose pas : elle nomme le billard, les echecs et
%  le trictrac sans les remplacer par le dhakon ou le macanan, qui
%  seraient les jeux d'ici.
%
%  SIXIEME REFUS DE LA COLONNE, DONC, ET IL EST DIFFERENT DES CINQ
%  AUTRES : ici ce n'est pas un mot javanais qu'on refuse pour une
%  chose etrangere, c'est une chose javanaise entiere — le jeu de la
%  cour, dhakon et macanan — qu'on refuse de substituer aux jeux
%  graves. Le refus porte sur la CHOSE et non plus seulement sur le
%  mot, et c'est la meme regle vue par l'autre bout.
%
%  L'HOTEL, LUI, EST NEERLANDAIS SANS PARTAGE : hotel, kamar, lift
%  l'ascenseur (9), kasir (14), rekening la facture (17), portir (18),
%  telepon (13), koper (31), omnibus (29), tilgram (34), montir le
%  mecanicien (52), setir. Une industrie nee au XIXe siecle porte le
%  vocabulaire du XIXe siecle, comme le chemin de fer du tableau 12 et
%  la ville du tableau 11. C'est le troisieme tableau de suite ou le
%  javanais emprunte plus qu'il ne garde, et les trois sont les trois
%  tableaux de la vie moderne.
%
%  MAIS « juru masak » (45) ET « pawon » (47) TIENNENT ENCORE. La
%  cuisine reste javanaise dans un hotel hollandais, comme elle
%  tenait au tableau 6 sous le kompor a gaz. C'est le dernier endroit
%  de la maison ou la langue ne cede pas.
% ===================================================================

%%K t13-noto-1 noto
\textit{\VUgras{(*) Bab rinciné kamar, delengen gambar No. 6.}}

%%K t13-apar-1 apar
\VUsaut{3.0mm}
\VUtitre{15.0pt}{60}{GAMBAR No. 13}
\VUsaut{1.6mm}
\VUpk{10.2pt}{40}{Hotel. --- Restoran.}
\VUsaut{2.0mm}
\VUpk{10.2pt}{40}{Warung Kopi.}
\VUsaut{3.4mm}

%%K t13-01-1 p
1. --- Sawise tekan Royan wingi sore, aku nginep ing
« \textit{Hotel Samodra} » kang wis disaranake marang aku dening
kancaku.

%%K t13-01-2 p
Aku diwenehi \VUgras{kamar} \textsuperscript{(2)} kang endah, ing
\VUgras{lantai} \textsuperscript{(1)} loro, panggonan aku turu kepenak
banget (*).

%%K t13-02-1 p
2. --- Esuk iki, metu saka kamarku, panggonan \VUgras{batur wadon}
\textsuperscript{(3)} wis nggawakake sarapan cilik, aku mlebu
\VUgras{kamar udud} \textsuperscript{(5)} kang uga dianggo \VUgras{kamar
maca} \textsuperscript{(4)} kanggo maca \VUgras{koran}
\textsuperscript{(6)} karo ngudud \VUgras{rokok} \textsuperscript{(8)}.
Sawise maca, aku mudhun nganggo \VUgras{lift} \textsuperscript{(9)} banjur
mlaku-mlaku ing kutha.

%%K t13-03-1 p
3. --- Sarehne aku lali takon marang \VUgras{juru tulis}
\textsuperscript{(12)} hotel bab jam pira panganan disuguhake ing
\VUgras{meja bareng} \textsuperscript{(11)}, aku bali luwih esuk lan
nganggo kesempatan iku kanggo adus ing \VUgras{kamar adus}
\textsuperscript{(10)} hotel. Banjur aku maneh menyang \VUgras{kantor
kasir} \textsuperscript{(15)} kanggo nelpon salah sijine kancaku. Nanging
\VUgras{telepon}e \textsuperscript{(13)} lagi dianggo ; weruh bingungku,
\VUgras{kasir wadon} \textsuperscript{(14)}, kang lagi nyathet
\VUgras{rekening} \textsuperscript{(17)} ing \VUgras{buku dhaptar tamu}
\textsuperscript{(16)}, njaluk marang \VUgras{portir}
\textsuperscript{(18)} supaya ngutus \VUgras{bocah hotel}
\textsuperscript{(19)} nindakake pesenku \VUgras{nganggo pit}
\textsuperscript{(20)}.

%%K t13-04-1 p
4. --- Aku matur nuwun marang dheweke banjur metu kanggo menehi
persenan marang \VUgras{kuli} \textsuperscript{(24)} kang wis nggawa saka
setasiun ing \VUgras{gerobag tangan}e \textsuperscript{(25)}
\VUgras{kothak topi}ku \textsuperscript{(26)}, \VUgras{koperku}
\textsuperscript{(27)} lan \VUgras{tas lelunganku} \textsuperscript{(28)}.
\VUgras{Tukang layang} \textsuperscript{(21)} kang lagi wae teka menehi
aku \VUgras{layang} \textsuperscript{(22)}.

%%K t13-05-1 p
5. --- Aku isih ngadeg sawetara wektu ing tlundhake lawang, cedhak
\VUgras{kothak layang} \textsuperscript{(23)}, nyawang mlaku-mlakune wong
ing dalan. \VUgras{Kondhektur} \textsuperscript{(30)} \VUgras{omnibus}
\textsuperscript{(29)} lagi ngunggahake ing kretane \VUgras{kopher}
\textsuperscript{(31)} \VUgras{wong lelungan} \textsuperscript{(32)} kang
lagi diterake \VUgras{sing duwe hotel} \textsuperscript{(33)}.
\VUgras{Tukang telegraf} \textsuperscript{(35)} enom kanthi ora
kesusu nggawa \VUgras{tilgram} \textsuperscript{(34)} kanggo tamu hotel ;
\VUgras{wong plesir} \textsuperscript{(36)} karo \VUgras{kamera}ne
\textsuperscript{(38)} nyampir, lagi maca \VUgras{buku panuntun}e
\textsuperscript{(37)}.

%%K t13-tit-1 sub
\VUsaut{4.0mm}
\VUpk{10.2pt}{40}{Jam sarapan awan.}
\VUsaut{3.0mm}

%%K t13-06-1 p
6. --- Ing ngarepe pager wesi, \VUgras{tukang komisi}
\textsuperscript{(39)} kang ora ambil pusing ngantuk ing antarane
\VUgras{cagak momotan}e \textsuperscript{(40)} lan \VUgras{kothak
semir}e \textsuperscript{(41)}, dene \VUgras{tukang umbah}
\textsuperscript{(42)} wadon enom kang nggawa \VUgras{kranjang}
\textsuperscript{(43)} sandhangan lawon ngguyu weruh polatan sarwa
resmine \VUgras{petugas meja} \textsuperscript{(44)} kang ngadeg cedhak
lawang.

%%K t13-07-1 p
7. --- Weruh \VUgras{juru masak} \textsuperscript{(45)} kang metu saka
\VUgras{pawon}e \textsuperscript{(47)} ing \VUgras{lantai ngisor lemah}
\textsuperscript{(48)} kanggo ngutus \VUgras{tukang isah-isah}
\textsuperscript{(46)} nindakake pesenan, aku eling yen jam sarapan awan
wis meh teka, banjur aku mlebu restoran, ngliwati latar panggonan
\VUgras{sopir} \textsuperscript{(50)} markir \VUgras{montor}e
\textsuperscript{(49)}, dene \VUgras{montir} \textsuperscript{(52)}
ndandani, miturut pituduhe \VUgras{tukang pit} \textsuperscript{(53)},
\VUgras{sepedha motor rodha telu} \textsuperscript{(51)} kang lagi wae
tabrakan.

%%K t13-08-1 p
8. --- Aku takon marang \VUgras{bocah hotel} \textsuperscript{(54)} enom
kang nggawa \VUgras{gelas bir} \textsuperscript{(55)}, apa meja bareng ana
ing sangisore emper, lan merga wangsulane iya, aku njupuk panggonan ing
salah sijine meja lan njaluk disuguhi panganan kang enak banget.

%%K t13-noto-2 noto
\textit{\VUgras{(*) Kanthi bantuan dhaptar panganan kang dilebokake ing
bausastra, guru bakal gampang ngowah-owahi menu ing ngisor iki.}}

%%K t13-tit-2 sub
\VUsaut{4.0mm}
\VUpk{10.2pt}{40}{Sarapan awan kang enak banget.}
\VUsaut{3.0mm}

%%K t13-09-1 p
9. --- Pelayan nggawakake menu (*) sarapan awan lan dhaptar anggur
marang aku. Aku milih lan mesen minangka \VUgras{panganan awal urang}
karo \VUgras{mentega,} lan minangka pambuka, \VUgras{babat cara Caen.}
Aku ora mangan \VUgras{iwak,} nanging njaluk saporsi \VUgras{pupu}
wedhus enom, lan minangka \VUgras{panganan sayur, bayem} karo krim.
Banjur, sarehne ora ana \VUgras{panganan selingan} kang cocog karo
seneng-senengku, aku njaluk digawakake warna-warna panganan panutup :
\VUgras{keju,} \VUgras{woh-wohan} lan \VUgras{biskuit.} Aku uga nambahi
\VUgras{anggur} Saint-Emilion kang enak banget, lan kanthi marem
banget aku ninggal meja sawise menehi \VUgras{pelayan}
\textsuperscript{(57)} \VUgras{persenan} \textsuperscript{(59)} kang lumayan.

%%K t13-10-1 p
10. --- Weruh aku wis siyap mangkat, \VUgras{sing ngurus}
\textsuperscript{(60)} ngajak aku munggah menyang balkon kanggo
nyawang panorama \VUgras{Samodra} \textsuperscript{(62)} kang endah
banget. Ing kadohan wong wiwit weruh \VUgras{prau mancing}
\textsuperscript{(63)}, \VUgras{kapal layar} \textsuperscript{(64)} lan
\VUgras{kapal paket} \textsuperscript{(65)} kang gedhe temen, kang silhuet
irenge nyembul ing \VUgras{mendhung} \textsuperscript{(67)} putih ing
\VUgras{wates langit} \textsuperscript{(66)}. Angin sumilir gawe obahe
\VUgras{gendera} \textsuperscript{(70)} kang dipacak ing dhuwure
\VUgras{papan jeneng} \textsuperscript{(69)} \VUgras{aula}
\textsuperscript{(68)}.

%%K t13-tit-3 sub
\VUsaut{3.0mm}
\VUpk{10.2pt}{40}{Panglipur.}
\VUsaut{3.0mm}

%%K t13-11-1 p
11. --- Tontonane narik banget, nganti aku mutusake sesuk munggah
menyang terase warung kopi karo nggawa \VUgras{teropong loro}
\textsuperscript{(71)} utawa \VUgras{teropong} \textsuperscript{(72)}.

%%K t13-12-1 p
12. --- Nalika ninggal balkon, aku menyang warung kopi hotel kanggo
ngombe \VUgras{omben-omben} \textsuperscript{(73)} karo main
\VUgras{biliar} \textsuperscript{(75)}. Tanpa mandheg aku ngliwati
ruwang kopi panggonan akeh tamu main \VUgras{catur}
\textsuperscript{(78)}, \VUgras{dham} \textsuperscript{(79)},
\VUgras{domino} \textsuperscript{(80)}, \VUgras{triktrak}
\textsuperscript{(81)} utawa \VUgras{kertu} \textsuperscript{(82)}, lan
aku langsung munggah menyang \VUgras{kamar biliar} \textsuperscript{(74)}.

%%K t13-13-1 p
13. --- Nalika permainane wis rampung, aku mudhun menyang lantai ngisor
lan njaluk marang pelayan \VUgras{amplop} \textsuperscript{(84)},
\VUgras{kertas} \textsuperscript{(83)}, \VUgras{prangko}
\textsuperscript{(85)}, \VUgras{pena} \textsuperscript{(87)} lan
\VUgras{mangsi} \textsuperscript{(86)} kanggo nulis layang.

%%K t13-13-2 p
Banjur aku ngundang \VUgras{kusir} \textsuperscript{(92)} kang
\VUgras{kreta}ne \textsuperscript{(88)} mandheg ing ngarepe warung kopi.
Aku takon marang dheweke bab rega miturut jam utawa sapisan mlaku, lan,
bareng kita wis sarujuk bab rega, dheweke mbukakake \VUgras{lawang
kreta} \textsuperscript{(89)} kanggo aku lan nglungguhake aku. Dheweke
munggah maneh ing \VUgras{jok}e \textsuperscript{(91)}, enggal-enggal
mangkatake jarane nganggo \VUgras{cemeti} \textsuperscript{(93)}, lan
kita mangkat mlaku-mlaku kang nyenengake.

\VUsaut{6.0mm}
\VUfilet{22mm}
